% Direct-authority diplomatic French for printed p.148.
\SourcePage{148}

\LayerTarget{tohoku:construction:composite-functor-spectral-sequences:p2}
\noindent $C(A)$. Il s'ensuit aussitôt que les suites spectrales définies par
$G$ et ce complexe $F(C(A))$ ne dépendent que de $A$, ce sont donc des
foncteurs spectraux cohomologiques sur $C$, ayant le même aboutissement,
appelés \textit{foncteurs spectraux du foncteur composé} $G\mathbin{\circ}F$.
Les formules données plus haut donnent immédiatement leurs termes initiaux :
\[
{}_{\mathrm I}E_2^{p,q}(A)=\bigl(R^p((R^qG)F)\bigr)(A),
\qquad
{}_{\mathrm{II}}E_2^{p,q}(A)=R^pG(R^qF(A)).
\]
De loin le cas le plus important pour l'obtention de suites spectrales non
triviales est celui où $F$ transforme objets injectifs en objets annulés par
les $R^qG$, $q\geq1$ (on appelle de tels objets $G$-acycliques), et où
$R^0G=G$ (i. e. $G$ est exact à gauche) : alors
${}_{\mathrm I}E_2^{p,q}=0$ si $q>0$, et se réduit à $R^p(GF)$ si $q=0$,
d'où résulte que l'aboutissement commun des deux suites spectrales s'identifie
au foncteur dérivé droit de $GF$. On obtient ainsi le

\LayerTarget{tohoku:thm:2.4.1}
\noindent\textsc{Théorème 2.4.1.} --- \textit{Soient $C,C',C''$ trois
catégories abéliennes, on suppose que tout objet de $C$ ou $C'$ est isomorphe
à un sous-truc d'un objet injectif. Considérons des foncteurs covariants
$F:C\longrightarrow C'$ et $G:C'\longrightarrow C''$, on suppose que $G$ est
exact à gauche et que $F$ transforme objets injectifs en objets $G$-acycliques
(i. e. annulés par les $R^qG$, $q>0$). Alors il existe une foncteur spectral
cohomologique sur $C$, à valeurs dans $C''$, aboutissant au foncteur dérivé à
droite $\mathfrak R(GF)$ de $GF$, (convenablement filtré), et dont le terme
initial est}
\LayerTarget{tohoku:eq:2.4.3}
\[
\tag{2.4.3}
E_2^{p,q}(A)=R^pG(R^qF(A)).
\]

\LayerTarget{tohoku:remark:2.4.1:1}
\noindent\textsc{Remarques 1.} La deuxième hypothèse sur le couple $(F,G)$
signifie que les foncteurs $(R^qG)F$ ($q>0$) sont effaçables (cf. 2.2), ou
encore que pour tout $A\in C$, on peut trouver un monomorphisme de $A$ dans un
$M$ tel que $(R^qG)(F(M))=0$ pour $q\geq1$ ; c'est sous cette forme que le plus
souvent nous vérifierons cette hypothèse.

\LayerTarget{tohoku:remark:2.4.1:2}
\noindent 2. On vérifie aussitôt que pour calculer la deuxième suite spectrale
d'un foncteur composé (i. e. celle dont il est question dans le théorème
2.4.1) il suffit de prendre une résolution $C(A)$ de $A$ par des $C^i$ qui
sont $F$-acycliques (et non nécessairement injectifs) et de prendre la
deuxième suite spectrale d'hyperhomologie du foncteur $G$ par rapport au
complexe $FC(A)$.

\LayerTarget{tohoku:remark:2.4.1:3}
\noindent 3. Notons deux cas particuliers importants où l'une des deux suites
spectrales d'hyperhomologie d'un foncteur $F$ par rapport à un complexe $K$
dégénère. \textit{Si $K$ est une résolution\textsuperscript{3)} d'un objet
$A$ de $C$, alors $\mathfrak R^nF(K)=R^nF(A)$}, donc l'objet gradué
$(R^nF(A))$ est l'aboutissement d'une suite spectrale de terme initial
${}_{\mathrm I}E_2^{p,q}(K)=H^p(R^qF(K))$. \textit{Si les $K^n$ sont
$F$-acycliques} (i. e. $R^mF(K^n)=0$ pour $m>0$) \textit{alors
$\mathfrak R^nF(K)=H^n(F(K))$}, donc l'objet gradué $(H^n(F(K)))$ est
l'aboutissement d'une suite spectrale de terme initial
${}_{\mathrm{II}}E_2^{p,q}(K)=R^pF(H^q(K))$. Combinant ces deux résultats, on
trouve donc : \textit{si $K$ est une résolution de $A$ par des objets
$F$-acycliques, on a $R^nF(A)=H^n(F(K))$}. Les isomorphismes ainsi obtenus sont
d'ailleurs aussi les morphismes déduits d'un homomorphisme de $K$ dans une
résolution injective de $A$ ; par suite, ils coïncident aussi, au signe près,
avec les homomorphismes définis par le ``iterated connecting homomorphism''
dans $K$ [6, V, 7].

